\begin{lemma}[Claim 4.1(a5): strong interchangeability]
Let \(n\) be a terminal PE round and \(L\) the active set.  Assume the
explicit lower-rung assertion (ii) needed here: for every \(t\in L\) with
\(t<n\), the colors \(\gamma_t,\delta_t\) are interchangeable on
\(T_{t+1}\) under \(\varphi_t\).  Let \(\sigma_n\) be a proper coloring
leaving precisely \(e_0\) uncolored and
\((T_n,D_n,\varphi_n)\)-stable.  Then for every \(t\in L\setminus\{n\}\),
\(\gamma_t\) and \(\delta_t\) are \(T_t\)-strongly interchangeable under
\(\sigma_n\), in the sense of \noderef{ETTStrongInterchangeable}.
\end{lemma}
\begin{proof}
Fix \(t\in L\) with \(t<n\).  Since the terminal round \(n\) is a PE, it
is not an RE.  Hence the finite set of indices (r>t) with
\(\Theta_r\ne\mathrm{RE}\) is nonempty.  Let \(\epsilon(t)\) be its least
member.  By construction
\[
 t<\epsilon(t)\leq n,qquad
 \Theta_{\epsilon(t)}\ne\mathrm{RE},
\]
and every strictly intermediate slot is an RE.  These are exactly the
first three witness clauses of
\noderef{ETTStrongInterchangeable}.

Inspect the PE certificate at (t) in \noderef{TashkinovSeriesStep}.
Its \noderef{ParallelExtensionColor} field makes \(\gamma_t\) missing at
the supporting vertex before the switch, while its selected connecting
edge supplies \(\delta_t\) there.  Thus the colors are distinct.  The
\noderef{TreeExit} field and the definition
\noderef{TwoColorComponentSwap} show that under \(\varphi_t\) the
\((\gamma_t,\delta_t)\)-component through \(v_t\) is a path meeting
\(T_t\) only at \(v_t\).  To see that the witnessed exit path is the
whole component, note that \(\gamma_t\) is missing at \(v_t\), while the
outside endpoint in \noderef{TreeExit} misses one of the two colors.
Both endpoints therefore have incidence-graph degree one; all degrees
are at most two by \noderef{TwoColorIncidenceDegreeAtMostTwo}.  Any
component vertex outside the witnessed path would attach to an internal
path vertex, creating degree at least three, or to an endpoint, creating
degree at least two.  Hence no such vertex exists.  The component remains
the same path after \noderef{TwoColorComponentSwap}, with its two colors
interchanged.

Now take \(v\in V(T_t)\setminus\{v_t\}\).  The stated lower-rung
interchangeability on \(T_{t+1}\), namely
\noderef{TreeInterchangeable}, permits at most one two-color path
component meeting that tree.  The path through \(v_t\) is already such a
component.  Therefore the component through \(v\) cannot be a path.
More explicitly, \noderef{TwoColorIncidenceDegreeAtMostTwo} bounds every
degree by two.  If the component had a vertex of degree at most one, the
finite traversal argument in \noderef{TwoColorComponentSpanningPath}
would make it a path component.  Thus every component vertex has degree
two, and following unused edges until the first repetition gives the
cycle required by \noderef{TwoColorComponentCycle}.  Distinct connected
components are disjoint, and hence this cycle avoids \(v_t\).

We claim that the cycle is contained in \(G[T_{\epsilon(t)}]\).  If not,
all intermediate slots are RE and hence preserve the coloring, by their
constructor equalities in \noderef{TashkinovSeriesStep}.  The same cycle
under \(\varphi_{\epsilon(t)-1}\) therefore meets the earlier tree \(T_t\)
but leaves \(T_{\epsilon(t)}\).  Traverse it from a vertex of \(T_t\),
take the first edge leaving \(T_{\epsilon(t)}\), and then the last edge
by which the traversal returns before ending at the chosen vertex of
\(T_t\).  The latter edge is an exact boundary
identity with one endpoint in and one endpoint out.  The segment from
that edge to the first subsequent vertex of \(T_t\) is simple after
deleting any closed subwalk, stays inside \(T_{\epsilon(t)}\), lies on
the same two-color cycle, and ends in the earlier tree.  The cycle,
leaving edge, endpoints, and this segment satisfy every field of
\noderef{RevisitingExtensionCertificate}: the exact earlier positive PE
index and its gamma/delta fields are those of \(t\); the cycle meets its
tree and has a vertex outside the current tree; and the chosen return
path is internally current.  Algorithm 3.1 gives RE priority.  In the
relational encoding this is the `available' constructor alternative in
\noderef{TashkinovSeriesStep}; consequently the slot
\(\epsilon(t)\) would have kind RE.  This contradicts the choice of
\(\epsilon(t)\).  Thus every vertex node of the cycle lies in
\(T_{\epsilon(t)}\).

It remains to transport this statement from \(\varphi_t\) to
\(\sigma_n\).  Both colors belong to the exact stored pair \(S_t\).  By
the definition \noderef{TashkinovStoredColors} and the monotonicity of
missing-plus-stored colors, they lie in the controlling set at every
stage from \(\epsilon(t)\) through \(n\), by
\noderef{HkuMissingStoredMonotone}.  The intermediate RE constructor
equalities identify
\(\varphi_t=\varphi_{\epsilon(t)-1}\).  Apply
\noderef{SameColorSeriesPreserved} from stage \(\epsilon(t)\) to every
exact edge identity on the contained cycle: its color is unchanged up to
\(\varphi_n\).  Terminal stability \noderef{StableColoring} then fixes
the same identities from \(\varphi_n\) to \(\sigma_n\).  Conversely,
symmetry of stability from \noderef{sc2}, together with properness,
prevents a new incident edge of either color at a cycle
vertex, so the full two-color component, not merely the displayed cycle,
is unchanged.  It remains a cycle, avoids \(v_t\), and is contained in
\(T_{\epsilon(t)}\).

The PE kind and exact support are supplied by active membership and the
chosen-vertex hypothesis.  Together with the witness
\(\epsilon(t)\) and the preceding component conclusions, these are all
clauses of \noderef{ETTStrongInterchangeable}.  This proves (a5).
\end{proof}
